% --- Slide 1 ---
\begin{frame}
\titlepage
\end{frame}

% --- Slide 2 ---
\begin{frame}{Nội dung trình bày}
\begin{enumerate}
    \item Nền tảng: 3D Gaussian là gì
    \item Chiếu phối cảnh \& Rasterizer khả vi
    \item Loss, Metric, Gradient \& Optimizer
    \item Adaptive Density Control (Clone/Split/Prune)
    \item Ba đòn bẩy tăng tốc của FastGS
    \item FastGS-lite trên Colab T4: pipeline \& kết quả
    \item Quan sát, giới hạn \& kết luận
\end{enumerate}
\end{frame}

% --- Slide 3 ---
\begin{frame}{Động lực: vì sao cần một biểu diễn cảnh mới?}
\begin{itemize}
    \item NeRF (Neural Radiance Fields): biểu diễn \textbf{implicit} bằng MLP, chất lượng cao nhưng render rất chậm (giây/khung hình).
    \item 3D Gaussian Splatting (3DGS): biểu diễn \textbf{explicit} bằng hàng trăm nghìn ``quả cầu'' Gaussian mờ trong không gian.
    \item Render bằng \textbf{rasterization khả vi} (differentiable rasterizer) thay vì ray-marching qua mạng nơ-ron.
    \item Kết quả: tốc độ render đạt hàng chục--hàng trăm FPS, gần thời gian thực.
\end{itemize}
\vspace{0.3cm}
\small
\centering
\begin{tabular}{lll}
\toprule
\textbf{Tiêu chí} & \textbf{NeRF} & \textbf{3D Gaussian Splatting} \\
\midrule
Biểu diễn & Implicit (MLP) & Explicit (tập Gaussian) \\
Render & Ray-marching, chậm & Rasterization, nhanh \\
Tốc độ & $\sim$giây/khung hình & Hàng chục--hàng trăm FPS \\
Huấn luyện & Backprop qua MLP & Backprop qua tham số Gaussian \\
\bottomrule
\end{tabular}
\end{frame}

% --- Slide 4 ---
\begin{frame}{Dự án Digital Twin GS / FastGS-lite}
\begin{itemize}
    \item Đây là bản \textbf{fork rút gọn} của FastGS, giữ nguyên thuật toán gốc.
    \item Mục tiêu: chạy vừa trên \textbf{1 GPU Colab T4 miễn phí} ($\sim$15GB VRAM) + 12.7GB RAM hệ thống.
    \item Phần cứng gốc của FastGS: GPU RTX 4090 (24GB VRAM).
\end{itemize}
\vspace{0.3cm}
\small
\centering
\begin{tabular}{lll}
\toprule
\textbf{Tiêu chí} & \textbf{Upstream FastGS} & \textbf{fastgs-lite} \\
\midrule
Phần cứng mục tiêu & RTX 4090, 24GB VRAM & Colab T4, $\sim$15GB VRAM \\
CUDA extension & Submodule ngoài & Vendor sẵn trong repo \\
Entry point & train.py & train.py + notebook có log \\
\bottomrule
\end{tabular}
\end{frame}

% --- Slide 5 ---
\begin{frame}{3D Gaussian là gì?}
\begin{columns}
\begin{column}{0.55\textwidth}
\footnotesize
Hàm mật độ Gaussian 3D:
\[
G(x) = \exp\left( -\frac{1}{2} (x-\mu)^T \Sigma^{-1} (x-\mu) \right)
\]
Mỗi Gaussian có 4 nhóm tham số học được:
\begin{itemize}
    \item Vị trí $\mu \in \mathbb{R}^3$
    \item Hiệp phương sai $\Sigma$ (hình dạng, hướng)
    \item Độ mờ (opacity) $\alpha$
    \item Màu qua hệ số Spherical Harmonics
\end{itemize}
\end{column}
\begin{column}{0.45\textwidth}
\centering
\includegraphics[width=0.95\linewidth]{ch2_gaussian_bump.png}
\captionof{figure}{\footnotesize Quả chuông mật độ Gaussian 3D}
\end{column}
\end{columns}
\end{frame}

% --- Slide 6 ---
\begin{frame}{Tham số hoá hiệp phương sai}
\begin{columns}
\begin{column}{0.55\textwidth}
\footnotesize
\begin{itemize}
    \item Không học trực tiếp $\Sigma$ để đảm bảo tính \textbf{positive semi-definite} trong lúc tối ưu.
    \item Thay vào đó học scale $s \in \mathbb{R}^3$ và quaternion xoay $q \in \mathbb{R}^4$:
    \[
    \Sigma = R S S^T R^T
    \]
    với $S = \mathrm{diag}(s)$, $R = R(q)$.
    \item Khởi tạo scale ban đầu dựa trên khoảng cách tới các điểm lân cận từ COLMAP.
\end{itemize}
\end{column}
\begin{column}{0.45\textwidth}
\centering
\includegraphics[width=0.78\linewidth]{ch1_scale_init.png}
\captionof{figure}{\footnotesize Khởi tạo scale từ khoảng cách lân cận}
\end{column}
\end{columns}
\end{frame}

% --- Slide 7 ---
\begin{frame}{Ellipsoid hiệp phương sai}
\begin{columns}
\begin{column}{0.55\textwidth}
\footnotesize
\begin{itemize}
    \item $\Sigma$ tương đương một \textbf{ellipsoid 3D} biểu diễn vùng ảnh hưởng của Gaussian.
    \item Trục ellipsoid = eigenvector của $\Sigma$.
    \item Độ dài trục = eigenvalue của $\Sigma$.
    \item Ellipsoid càng dẹt/kéo dài $\Rightarrow$ Gaussian càng bất đẳng hướng.
\end{itemize}
\end{column}
\begin{column}{0.45\textwidth}
\centering
\includegraphics[width=0.95\linewidth]{ch2_covariance_ellipsoid.png}
\captionof{figure}{\footnotesize Ellipsoid ứng với $\Sigma$}
\end{column}
\end{columns}
\end{frame}

% --- Slide 8 ---
\begin{frame}{Màu sắc qua Spherical Harmonics (SH)}
\begin{columns}
\begin{column}{0.55\textwidth}
\footnotesize
\begin{itemize}
    \item Mỗi Gaussian không lưu 1 màu RGB cố định, mà lưu hệ số \textbf{SH}.
    \item Biểu diễn màu \textbf{phụ thuộc góc nhìn} (view-dependent), đặc tả hiệu ứng phản xạ/specular.
    \item Hệ số SH bậc 0 (DC) khởi tạo trực tiếp từ màu RGB quan sát được:
    \[
    f_{dc} = \mathrm{RGB2SH}(\text{color})
    \]
\end{itemize}
\end{column}
\begin{column}{0.45\textwidth}
\centering
\includegraphics[width=0.95\linewidth]{ch1_rgb2sh.png}
\captionof{figure}{\footnotesize Chuyển đổi RGB $\to$ hệ số SH bậc 0}
\end{column}
\end{columns}
\end{frame}

% --- Slide 9 ---
\begin{frame}{Cơ sở SH và lịch tăng bậc}
\footnotesize
\begin{columns}
\begin{column}{0.5\textwidth}
\centering
\includegraphics[width=0.9\linewidth]{ch2_sh_sphere.png}
\captionof{figure}{\footnotesize Các hàm cơ sở SH bậc thấp trên mặt cầu}
\end{column}
\begin{column}{0.5\textwidth}
\centering
\includegraphics[width=0.9\linewidth]{ch2_sh_degree_schedule.png}
\captionof{figure}{\footnotesize Lịch tăng dần bậc SH khi huấn luyện}
\end{column}
\end{columns}
\vspace{0.2cm}
\small
Mỗi vài nghìn iteration tăng thêm 1 bậc SH, tối đa bậc 3 (16 hệ số/kênh màu).
\end{frame}

% --- Slide 10 ---
\begin{frame}{Khởi tạo từ SfM/COLMAP}
\begin{columns}
\begin{column}{0.55\textwidth}
\footnotesize
\begin{itemize}
    \item Gaussian ban đầu được khởi tạo tại các điểm \textbf{point-cloud thưa} do COLMAP tái tạo (Structure-from-Motion).
    \item \textbf{Scene extent}: đường chéo bounding box của các camera trong cảnh.
    \item Scene extent dùng làm tham chiếu cho ngưỡng phân tách Gaussian (\texttt{-{}-dense}) ở Adaptive Density Control (phần sau).
\end{itemize}
\end{column}
\begin{column}{0.45\textwidth}
\centering
\includegraphics[width=0.95\linewidth]{ch1_extent.png}
\captionof{figure}{\footnotesize Scene extent từ bounding box camera}
\end{column}
\end{columns}
\end{frame}
